% preprints/jphq_05_kpi_limits/sections/02_formal_setting.tex

\section{Formal Setting}
\label{sec:formal-setting}

Throughout this preprint we assume the Ontology of Continua (OC) together with
the frozen executable theory specification \texttt{ETS.K\_EA.v1.1}.
An enterprise is modeled as a continuum \(K_{EA}\) with state space
\(\Omega(K_{EA})\), boundary \(\partial\Omega(K_{EA})\), axes \(A\), flows \(J\),
thresholds \(\Theta\), cycles \(C\), and a continuum measure \(k_{EA}\).
No primitives beyond those defined in OC and ETS are introduced or assumed.

\subsection{ETS-anchored observability structures}

In \texttt{ETS.K\_EA.v1.1}, observability is not treated as an external modeling
assumption but as an internal structural component of the enterprise continuum.
In particular, the executable specification includes:

\begin{itemize}[leftmargin=2em]
  \item an information-axis component \emph{observability}, located within the
        information/semantics axis block;
  \item a factor \(S_{\mathrm{obs}}\) in the factorization of the continuum
        measure \(k_{EA}\);
  \item an explicit \emph{ObservabilityInvariant} listed among the system
        invariants;
  \item and a phase-logic condition \texttt{observability\_lost} acting as a STOP
        trigger.
\end{itemize}

Observation is therefore defined only on the OC/ETS-admissible observable domain
\[
\Omega_{\mathrm{obs}} \;\subseteq\; \Omega(K_{EA}),
\]
consisting of all enterprise states for which \(S_{\mathrm{obs}}\) is defined
and the ObservabilityInvariant holds.
When the ETS phase condition \texttt{observability\_lost} applies, KPI readout is
outside scope and no observation map is assumed to be defined.

\subsection{Observation as structural quotienting}

Fix a finite KPI family \(K=(\kappa_1,\ldots,\kappa_m)\) admissible under OC/ETS.
Each KPI \(\kappa_i\) is treated as a readout map
\(\kappa_i:\Omega_{\mathrm{obs}}\to\mathbb{R}\).
Together they define the KPI readout map
\[
\pi_K:\Omega_{\mathrm{obs}}\to\Sigma_{\mathrm{obs}}\subseteq\mathbb{R}^m,
\qquad
\pi_K(\omega)
=
\bigl(\kappa_1(\omega),\ldots,\kappa_m(\omega)\bigr).
\]

As in classical observation theory, the map \(\pi_K\) induces an equivalence
relation \(\sim_K\) on \(\Omega_{\mathrm{obs}}\) defined by
\[
\omega_1 \sim_K \omega_2
\quad\Longleftrightarrow\quad
\pi_K(\omega_1)=\pi_K(\omega_2).
\]
Accordingly, KPI-based observation partitions \(\Omega_{\mathrm{obs}}\) into
fibers (preimages) of \(\pi_K\).
The observation space \(\Sigma_{\mathrm{obs}}\) is therefore identified with the
quotient
\[
\Sigma_{\mathrm{obs}} \;\cong\; \Omega_{\mathrm{obs}} / \!\sim_K ,
\]
which represents the maximal information preserved by the KPI readout
\cite{kalman1960,ljung1999}.

All subsequent limits established in this preprint are statements about this
quotienting.
Identifiability of enterprise states requires injectivity of \(\pi_K\) on a
domain of interest.
In contrast, KPI-based observation generically collapses structurally distinct
states in \(\Omega_{\mathrm{obs}}\) into the same equivalence class.
